\chapter{Il calcolatore}
\label{chap:calcolatore}

Per imparare a guidare non serve saper costruire un motore, ma aiuta molto
sapere che cosa fanno il volante, i pedali e il cambio. Allo stesso modo, per
programmare non serve conoscere l'elettronica di un computer, ma bisogna
avere un'idea chiara di come è fatto e di come esegue un programma: molte
regole del C, che altrimenti sembrerebbero capricci, diventano ovvie. Questo
capitolo costruisce quell'idea, partendo da che cosa significa
\textit{programmare}.

\section{Che cos'è la programmazione}

Il corso si apre con la definizione del dizionario Treccani:

\begin{center}
\begin{minipage}{0.86\linewidth}
\itshape “In informatica, la stesura di un programma in una forma
interpretabile dal calcolatore elettronico (o direttamente o dopo una
rielaborazione opportuna), compiuta dal programmatore sulla base delle
specifiche di programma fornite dall'analista; la programmazione comprende
la definizione del problema da risolvere, la preparazione del diagramma di
flusso del programma, la sua codifica (cioè la trascrizione, in un
linguaggio opportuno, delle varie istruzioni), la messa a punto e la
preparazione di una documentazione a uso dell'operatore.”
\end{minipage}
\end{center}

La definizione contiene già tutto il percorso del corso. Scrivere codice
(la \textit{codifica}) è solo una delle fasi, e non la prima: prima bisogna
capire il problema e progettare la soluzione, dopo bisogna verificare che
funzioni (la \textit{messa a punto}) e spiegarla (figura~\ref{fig:fasi}).
La parentesi “o direttamente o dopo una rielaborazione opportuna” è
importante: il programma che scriviamo non è ancora ciò che il computer
esegue, va \textit{rielaborato}, cioè tradotto. È il compito del
compilatore incontrato nel capitolo precedente.

\begin{figure}[H]
  \centering
  \begin{tikzpicture}[font=\Lato\scriptsize,
      fase/.style={wtnode, font=\Lato\scriptsize, text width=2.55cm,
        minimum height=1.5cm}]
    \node[fase] (a) at (0,0) {\textbf{definire il problema}\\[2pt]che cosa si deve ottenere, a partire da quali dati};
    \node[fase] (b) at (3.6,0) {\textbf{progettare}\\[2pt]l'algoritmo, per esempio con un diagramma di flusso};
    \node[fase, draw=accent!60!black, fill=accent!8!white] (c) at (7.2,0) {\textbf{codificare}\\[2pt]scrivere il programma in un linguaggio (il C)};
    \node[fase] (d) at (10.8,0) {\textbf{mettere a punto}\\[2pt]provare, trovare e correggere gli errori};
    \node[fase] (e) at (14.4,0) {\textbf{documentare}\\[2pt]spiegare come si usa e come funziona};
    \foreach \x/\y in {a/b, b/c, c/d, d/e} \draw[wtarrow] (\x) -- (\y);
    \draw[wtarrowacc] (d.south) -- ++(0,-0.4) -| node[wtlab, pos=0.25] {gli errori trovati riportano indietro} (b.south);
  \end{tikzpicture}
  \caption{Le fasi della programmazione secondo la definizione: scrivere il
           codice è solo una di esse.}
  \label{fig:fasi}
\end{figure}

\info{Come si impara a programmare}{
  A lezione il docente avverte che il corso parte da zero e all'inizio va
  piano, ma che velocità e difficoltà crescono presto, e che entro il primo
  mese quasi tutti incontrano difficoltà serie, soprattutto negli esercizi.
  È normale e si supera, con alcune abitudini:
  \begin{itemize}
    \item non dare per scontato di saper già programmare, e non saltare le
          lezioni “facili”;
    \item studiare con costanza: la programmazione si accumula, ogni
          argomento usa i precedenti;
    \item \textbf{scrivere ed eseguire} ogni esempio, poi modificarlo per
          vedere che cosa succede: leggere un programma non basta, come non
          basta guardare qualcuno nuotare;
    \item chiedere appena qualcosa non è chiaro (“non siete gli unici”);
    \item prendere appunti e fare sintesi: le slide sono, per loro natura,
          sintetiche; questa dispensa prova a riempire gli spazi, ma non
          sostituisce le lezioni.
  \end{itemize}}

\section{Lo schema di un calcolatore}

Un computer moderno è un oggetto complicatissimo, ma il suo funzionamento si
capisce con uno schema di tre soli blocchi (figura~\ref{fig:vonneumann}):

\begin{itemize}
  \item la \term{CPU} (\textit{Central Processing Unit}, il
        \textit{processore}), che esegue le istruzioni;
  \item la \term{memoria centrale} (o \textit{principale}, la RAM), che
        contiene le istruzioni da eseguire e i dati su cui lavorano;
  \item i \term{dispositivi di ingresso e uscita} (\textit{I/O},
        \textit{Input/Output}): tastiera, mouse, schermo, disco, rete\dots,
        collegati attraverso delle \textit{interfacce}.
\end{itemize}

I blocchi comunicano attraverso il \term{bus}, un insieme di collegamenti
condiviso su cui viaggiano dati, indirizzi e comandi. Fisicamente, è la
\textit{scheda madre} (\textit{motherboard}) a collegare fra loro processore,
memoria e periferiche.

\begin{figure}[H]
  \centering
  \begin{tikzpicture}[font=\Lato\scriptsize,
      dev/.style={wtnodeplain, minimum width=2.3cm, minimum height=0.75cm},
      intf/.style={draw=green-gradient-6, fill=green-gradient-1!45!white,
        line width=0.7pt, minimum width=2.3cm, minimum height=0.45cm,
        font=\Lato\tiny},
      blk/.style={wtnode, minimum width=3.4cm, minimum height=1.2cm}]
    \foreach \x/\n in {0/tastiera, 3/mouse, 6/schermo, 9/disco (memoria secondaria)} {
      \node[dev, text width=2.1cm] (d\x) at (\x,2.2) {\n};
      \node[intf] (i\x) at (\x,1.1) {interfaccia di I/O};
      \draw[wtarrow, <->] (d\x) -- (i\x);
      \draw[wtarrow, <->] (i\x) -- (\x,0.22);
    }
    \fill[orange-gradient-4] (-1.4,0.2) rectangle (10.4,-0.2);
    \node[font=\Lato\scriptsize\bfseries, text=orange-gradient-6!35!black] at (4.5,0) {bus};
    \node[blk, draw=red-gradient-6!70!black, fill=red-gradient-1!25!white] (mem) at (2.2,-1.6) {\textbf{memoria centrale}\\istruzioni e dati};
    \node[blk, draw=accent!60!black, fill=accent!8!white] (cpu) at (6.8,-1.6) {\textbf{CPU}\\esegue le istruzioni};
    \draw[wtarrow, <->] (mem) -- (2.2,-0.22);
    \draw[wtarrow, <->] (cpu) -- (6.8,-0.22);
  \end{tikzpicture}
  \caption{Schema semplificato di un calcolatore (architettura di von
           Neumann): CPU, memoria e dispositivi di I/O collegati dal bus.}
  \label{fig:vonneumann}
\end{figure}

Questo schema porta il nome di \textit{John von Neumann}, che nel 1945 ne
descrisse l'idea centrale: le \textbf{istruzioni} del programma stanno
nella \textbf{stessa memoria} dei \textbf{dati}. Un programma non è cablato
nei circuiti: è una sequenza di numeri in memoria, che si può caricare,
sostituire, copiare. È questo che rende un computer una macchina
\textit{generale}, capace di fare cose diversissime a seconda del programma
caricato.

\section{La CPU e il linguaggio macchina}

\subsection{Che cosa sa fare un processore}

Tutte le elaborazioni di un calcolatore sono svolte dalla CPU, oppure da
altri componenti su suo comando. Sorprendentemente, un processore sa
compiere solo operazioni \textbf{molto semplici}:

\begin{itemize}
  \item leggere, scrivere o copiare una o più celle di memoria;
  \item sommare, sottrarre, moltiplicare, dividere il contenuto di celle di
        memoria;
  \item leggere e scrivere in zone di memoria “speciali” per pilotare i
        dispositivi di I/O (per esempio la scheda video);
  \item poche altre operazioni elementari, come confrontare due numeri e
        \textbf{saltare} a un'altra istruzione a seconda del risultato.
\end{itemize}

Il segreto è la velocità: un processore moderno esegue \textbf{miliardi} di
queste operazioni al secondo. Ogni cosa che un computer fa, dal mostrare
una pagina web al riprodurre un video, è scomposta in una lunghissima
sequenza di passi elementari.

\subsection{Il linguaggio macchina}

Ogni processore ha un proprio insieme di istruzioni, e ogni istruzione è
identificata da una certa configurazione di bit (di zeri e uni). L'insieme
delle istruzioni di un processore si chiama \term{linguaggio macchina} di
quel processore. Processori di famiglie diverse hanno linguaggi macchina
diversi: un programma tradotto per i processori Intel e AMD dei PC (famiglia
\textit{x86-64}) non gira su un Mac con processore Apple (famiglia
\textit{ARM}), e viceversa.

Per eseguire un programma, allora, bisogna fare due cose: mettere in
memoria la sequenza di configurazioni di bit che rappresentano le sue
istruzioni, e dire al processore a quale indirizzo si trova la prima. Da lì
in poi il processore ripete sempre lo stesso ciclo
(figura~\ref{fig:ciclo-cpu}): \textbf{preleva} l'istruzione all'indirizzo
indicato da un suo registro speciale, il \term{program counter};
\textbf{decodifica} l'istruzione, cioè capisce che cosa chiede;
la \textbf{esegue}; e passa all'istruzione successiva, a meno che
l'istruzione appena eseguita non fosse un \textit{salto}, che cambia il
program counter.

\begin{figure}[H]
  \centering
  \begin{tikzpicture}[font=\Lato\scriptsize]
    % memoria
    \memoria{m}{(0,0)}{1000/MOV \ldots/, 1003/ADD \ldots/, 1005/JMP 1000/, 1008/HLT/, \ldots/\ldots/}
    \node[mnota, anchor=south] at ([yshift=2pt]m-1.north) {\textbf{memoria}};
    \node[mcellam, fill=yellow-gradient-1!70!white] at (m-2) {ADD \ldots};
    % CPU
    \node[wtnode, draw=accent!60!black, fill=accent!6!white, minimum width=4.6cm,
          minimum height=2.75cm, minimum width=5.3cm, anchor=north west] (cpu) at (3.0,0.55) {};
    \node[font=\Lato\scriptsize\bfseries, anchor=north west] at ([shift={(5pt,-4pt)}]cpu.north west) {CPU};
    \node[mvar] (pc) at (4.55,-0.49) {1003};
    \node[mnome, anchor=west] at ([xshift=3pt]pc.east) {program counter};
    \node[mvar] (ax) at (4.55,-1.2) {42};
    \node[mnome, anchor=west] at ([xshift=3pt]ax.east) {registro AX};
    \node[mvar] (bx) at (4.55,-1.8) {8};
    \node[mnome, anchor=west] at ([xshift=3pt]bx.east) {registro BX};
    \draw[mrif] (pc.west) -- ([xshift=8pt]m-2.east);
    \node[mpunto] at (pc.west) {};
    % ciclo
    \begin{scope}[shift={(10.9,-0.95)}]
      \node[wtnodeplain] (f) at (90:1.25) {1. preleva};
      \node[wtnodeplain] (d) at (-30:1.25) {2. decodifica};
      \node[wtnodeplain] (e) at (210:1.25) {3. esegue};
      \draw[wtarrow] (f) to[bend left=30] (d);
      \draw[wtarrow] (d) to[bend left=30] (e);
      \draw[wtarrow] (e) to[bend left=30] (f);
    \end{scope}
  \end{tikzpicture}
  \caption{Il program counter indica l'indirizzo della prossima istruzione;
           la CPU la preleva, la decodifica e la esegue, poi ricomincia. I
           registri sono piccole celle di memoria interne alla CPU.}
  \label{fig:ciclo-cpu}
\end{figure}

\section{L'assembly}

Il linguaggio macchina è una sequenza di bit, praticamente illeggibile per
una persona. Il primo passo per renderlo comprensibile è l'\term{assembly}:
un linguaggio quasi identico al linguaggio macchina, in cui però ogni codice
operativo è sostituito da un'abbreviazione (un \textit{codice mnemonico}):
\code{MOV} per \textit{move} (copia), \code{ADD} per \textit{add} (somma),
\code{JMP} per \textit{jump} (salta). Un programma, l'\term{assembler},
traduce l'assembly in linguaggio macchina, un'istruzione alla volta.

\subsection{Un esempio: contare i bit a 1}

Il programma seguente, visto a lezione, conta quanti bit valgono 1 in un
dato di 16 bit. A sinistra di ogni istruzione c'è il suo indirizzo in
memoria (in esadecimale, cioè in base 16: per ora basta considerarlo un
“numero civico”), a destra la spiegazione.

\begin{center}
\renewcommand{\arraystretch}{1.25}
{\FiraCodeNL\scriptsize
\begin{tabular}{@{}l l l >{\Lato\scriptsize}l@{}}
  \toprule
  \normalfont\Lato\bfseries indirizzo & \normalfont\Lato\bfseries istruzione & & \normalfont\Lato\bfseries significato \\
  \midrule
  3F0D0100 & 00000000 00000000 & & il dato da elaborare (16 bit) \\
  3F0D0102 & 00000000 & & qui andrà il risultato (8 bit) \\
  \midrule
  40B1A200 & MOV & \$0, \%AL         & azzera il contatore, il registro AL \\
  40B1A203 & MOV & (3F0D0100), \%DX  & copia nel registro DX il dato da elaborare \\
  40B1A207 & CMP & \%DX, \$0         & confronta DX con 0\dots \\
  40B1A211 & JE  & 40B1A232          & \dots e se sono uguali salta alla fine \\
  40B1A215 & SHL & \%DX              & sposta i bit di DX di un posto a sinistra: \\
           &     &                   & il bit che “esce” finisce nel \textit{carry flag} CF \\
  40B1A217 & JC  & 40B1A225          & se CF vale 1 (è uscito un 1) salta all'incremento \\
  40B1A221 & JMP & 40B1A207          & altrimenti torna al confronto \\
  40B1A225 & ADD & \$1, \%AL         & incrementa il contatore AL \\
  40B1A228 & JMP & 40B1A207          & torna al confronto \\
  40B1A232 & MOV & \%AL, (3F0D0102)  & scrive il risultato in memoria \\
  40B1A236 & HLT &                   & fine (\textit{halt}) \\
  \bottomrule
\end{tabular}}
\end{center}

L'idea è questa: a ogni giro si fa “scivolare” il dato di un posto a
sinistra, e si guarda il bit che cade fuori dal bordo; se è un 1, si
incrementa il contatore. Quando nel dato non resta nessun 1 (cioè vale 0),
si è finito. Seguiamo l'esecuzione con un dato che contiene due bit a 1,
\code{1010000000000000}:

\begin{center}
\renewcommand{\arraystretch}{1.3}
\begin{tabular}{@{}c>{\FiraCodeNL\small}c>{\FiraCodeNL\small}c>{\FiraCodeNL\small}c>{\FiraCodeNL\small}c>{\raggedright\arraybackslash}p{4.4cm}@{}}
  \toprule
  \textbf{giro} & \normalfont\textbf{DX prima} & \normalfont\textbf{bit uscito (CF)} & \normalfont\textbf{DX dopo} & \normalfont\textbf{AL} & \textbf{che cosa succede} \\
  \midrule
  1 & 1010000000000000 & 1 & 0100000000000000 & 1 & è uscito un 1: \code{JC} salta all'\code{ADD} \\
  2 & 0100000000000000 & 0 & 1000000000000000 & 1 & è uscito uno 0: \code{JMP} torna al confronto \\
  3 & 1000000000000000 & 1 & 0000000000000000 & 2 & è uscito un 1: si incrementa \\
  4 & 0000000000000000 & & & 2 & DX vale 0: \code{JE} salta alla fine \\
  \bottomrule
\end{tabular}
\end{center}

Il risultato, 2, viene scritto in memoria e il programma si ferma. Si noti
che in assembly non esiste un'istruzione “ripeti finché”: la ripetizione
si ottiene \textbf{saltando indietro} (\code{JMP 40B1A207}) e la scelta
\textbf{saltando avanti} se una condizione è vera (\code{JE}, \code{JC}).

\subsection{I limiti dell'assembly}

Con l'assembly si può scrivere qualunque programma, anche un sistema
operativo o un'interfaccia grafica, ma è un lavoro enorme: le istruzioni
sono di \textbf{basso livello}, cioè vicine alla macchina e lontane dal
modo in cui ragioniamo. Per “contare gli uni” abbiamo dovuto pensare in
termini di registri, flag e indirizzi di salto; e il programma funziona
solo sul processore per cui è scritto.

\section{I livelli di astrazione}

Servono istruzioni più vicine al modo di ragionare umano. La soluzione è
l'\term{astrazione}: costruire un linguaggio che “nasconde” la macchina
sottostante e permette di esprimere gli algoritmi in modo simbolico. Sono i
\term{linguaggi di alto livello}, come il C, che hanno due proprietà
fondamentali:

\begin{itemize}
  \item sono più vicini al linguaggio umano (e matematico) che al
        linguaggio macchina;
  \item sono \textbf{generali}: non dipendono dal processore. Lo stesso
        programma C si traduce per un PC, per un Mac o per un telefono.
\end{itemize}

La figura~\ref{fig:livelli} mostra la stessa istruzione ai tre livelli. È
la traduzione reale prodotta da gcc su un PC (processore x86-64) per
l'istruzione C \code{c = a + b;}, in cui \code{a}, \code{b} e \code{c} sono
tre variabili intere.

\begin{figure}[H]
  \centering
  \begin{tikzpicture}[font=\Lato\scriptsize,
      liv/.style={rounded corners=4pt, line width=0.8pt, minimum width=17.2cm,
        inner sep=6pt, align=left, anchor=north west},
      tit/.style={font=\Lato\scriptsize\bfseries, anchor=north west}]
    \node[liv, draw=green-gradient-6, fill=green-gradient-1!35!white] (alto) at (0,0)
      {\begin{tabular}{@{}p{3.5cm}l@{}}
         \textbf{linguaggio di alto livello} & {\FiraCodeNL\small c = a + b;} \\
         \textit{si ragiona sul problema} & \\
       \end{tabular}};
    \node[liv, draw=primary!65!black, fill=primary!6!white] (asm) at ([yshift=-0.35cm]alto.south west)
      {\begin{tabular}{@{}p{3.5cm}>{\FiraCodeNL\small}l>{\Lato\scriptsize}l@{}}
         \textbf{assembly} & mov edx, [rbp-0xc] & copia \code{a} nel registro edx \\
         \textit{si ragiona sulla macchina:} & mov eax, [rbp-0x8] & copia \code{b} nel registro eax \\
         \textit{registri e indirizzi} & add eax, edx & eax = eax + edx \\
         & mov [rbp-0x4], eax & copia il risultato in \code{c} \\
       \end{tabular}};
    \node[liv, draw=accent!60!black, fill=accent!6!white] (mac) at ([yshift=-0.35cm]asm.south west)
      {\begin{tabular}{@{}p{3.5cm}>{\FiraCodeNL\small}l@{}}
         \textbf{linguaggio macchina} & 10001011 01010101 11110100 \\
         \textit{bit, cioè numeri} & 10001011 01000101 11111000 \\
         & 00000001 11010000 \\
         & 10001001 01000101 11111100 \\
       \end{tabular}};
    \draw[wtarrowacc, line width=1.2pt] ([xshift=-0.6cm, yshift=-0.2cm]alto.north east) -- node[wtlab, left=4pt, fill=none, align=right] {traduzione\\(compilatore)} ([xshift=-0.6cm, yshift=0.2cm]mac.south east);
  \end{tikzpicture}
  \caption{La stessa istruzione a tre livelli di astrazione. Un'istruzione
           di alto livello diventa più istruzioni macchina; ogni istruzione
           assembly corrisponde a un'istruzione macchina.}
  \label{fig:livelli}
\end{figure}

Ogni livello richiede di conoscere cose diverse:

\begin{itemize}
  \item per il \textbf{linguaggio macchina}, come sono rappresentate le
        informazioni in bit;
  \item per l'\textbf{assembly}, le caratteristiche della macchina:
        registri, insieme di istruzioni, dimensione dei dati;
  \item per un \textbf{linguaggio di alto livello}, solo il linguaggio: ci
        si astrae dai dettagli della macchina.
\end{itemize}

Il prezzo da pagare è che un programma di alto livello non può essere
eseguito direttamente: va \textbf{tradotto} in linguaggio macchina. È il
compito del \term{compilatore}, l'argomento del
capitolo~\ref{chap:dal-problema}.

\subsection{Algoritmo e programma}

Due parole che useremo continuamente vanno distinte con cura.

\dfn{Algoritmo}{
  Una sequenza di passi elementari che risolve un problema. La sequenza
  deve essere \textbf{finita}, \textbf{non ambigua},
  \textbf{deterministica} e \textbf{generale}.}

\dfn{Programma}{
  La formalizzazione di un algoritmo in un linguaggio di programmazione.}

Le quattro proprietà dell'algoritmo non sono pignolerie: ciascuna esclude un
modo concreto in cui un procedimento può fallire. Prendiamo una ricetta di
cucina, che è quasi un algoritmo:

\begin{itemize}
  \item \textbf{finita}: deve terminare. “Mescola finché non è pronto” è
        accettabile solo se prima o poi “è pronto”;
  \item \textbf{non ambigua}: ogni passo deve avere un solo significato.
        “Sale quanto basta” va bene per un cuoco, non per un computer, che
        non ha buon senso: bisogna dire quanti grammi;
  \item \textbf{deterministica}: in ogni momento il passo successivo è
        determinato; con gli stessi dati si ottiene sempre lo stesso
        risultato;
  \item \textbf{generale}: risolve una \textit{classe} di problemi, non un
        solo caso. Una ricetta “per 4 persone” è meno generale di una “per
        $n$ persone”; un algoritmo che calcola l'area di un cerchio di
        raggio 2 è inutile, ne serve uno che funzioni per qualunque raggio.
\end{itemize}

\section{La memoria}

\subsection{Celle e indirizzi}

La memoria centrale è il contenitore dei dati su cui lavora il processore.
Si può immaginare come una lunghissima fila di caselle consecutive, le
\term{locazioni} (o \textit{celle}) di memoria. Ogni cella:

\begin{itemize}
  \item è l'unità minima di memorizzazione: si legge e si scrive una cella
        intera, non un pezzo;
  \item contiene tipicamente \textbf{8 bit}, cioè un \term{byte};
  \item è identificata da un numero, il suo \term{indirizzo}, come le case
        di una strada dal numero civico.
\end{itemize}

\begin{affianco}[0.5][c]
\begin{immagine}
\begin{tikzpicture}
  \memoria{m}{(0,0)}{1000/00000101/, 1001/01000001/, 1002/11111111/, 1003/00000000/, 1004/\dots/}
  \node[mnota, anchor=west] at ([xshift=12pt]m-1.east) {il numero 5};
  \node[mnota, anchor=west] at ([xshift=12pt]m-2.east) {il numero 65, o la lettera A};
  \node[mnota, anchor=west] at ([xshift=12pt]m-3.east) {il numero 255};
  \node[mnota, anchor=south] at ([yshift=3pt]m-1.north) {\textbf{contenuto} (8 bit)};
  \node[mnota, anchor=south east] at ([yshift=3pt, xshift=-4pt]m-1.north west) {\textbf{indirizzo}};
\end{tikzpicture}
\end{immagine}
\risultato
Un \term{bit} (\textit{binary digit}, cifra binaria) è la più piccola unità
di informazione: vale 0 oppure 1. Elettronicamente corrisponde a qualcosa
che ha due stati, come un interruttore acceso o spento. Con un bit solo si
distinguono due situazioni; mettendone insieme 8 se ne distinguono molte di
più: vediamo quante, e come si usano per rappresentare i numeri.
\end{affianco}

\subsection{La notazione binaria}

Siamo abituati alla notazione \textbf{decimale}: dieci cifre, da 0 a 9, e
un sistema \textbf{posizionale}, in cui il valore di una cifra dipende dalla
posizione. In \code{352} il 3 vale trecento perché sta nella posizione delle
centinaia:
\[
  352 = 3 \cdot 100 + 5 \cdot 10 + 2 \cdot 1 = 3 \cdot 10^2 + 5 \cdot 10^1 + 2 \cdot 10^0 .
\]
La notazione \textbf{binaria} funziona allo stesso modo, ma con due sole
cifre, 0 e 1, e con le potenze di \textbf{2} al posto delle potenze di 10
(figura~\ref{fig:binario}).

\begin{figure}[H]
  \centering
  \begin{tikzpicture}[font=\FiraCodeNL\small]
    \foreach \p [count=\i from 0] in {128,64,32,16,8,4,2,1} {
      \node[mcella, minimum width=1.25cm, minimum height=0.7cm, font=\FiraCodeNL\normalsize]
        (b\i) at (\i*1.25,0) {};
      \node[font=\FiraCodeNL\scriptsize, text=neutral!40!black] at (\i*1.25,0.62) {\p};
      \pgfmathtruncatemacro{\e}{7-\i}
      \node[font=\Lato\tiny, text=neutral!50!black] at (\i*1.25,1.0) {$2^{\e}$};
    }
    \foreach \val [count=\i from 0] in {0,0,1,0,1,1,0,1} {
      \ifnum\val=1
        \node[font=\FiraCodeNL\normalsize\bfseries, text=accent!80!black] at (b\i) {1};
      \else
        \node[font=\FiraCodeNL\normalsize, text=neutral!55!black] at (b\i) {0};
      \fi
    }
    \node[anchor=west, font=\Lato\small] at (10.1,0)
      {$= 32 + 8 + 4 + 1 = \mathbf{45}$};
    \node[font=\Lato\tiny, anchor=east, text=neutral!50!black] at (-0.75,1.0) {potenza};
    \node[font=\Lato\tiny, anchor=east, text=neutral!50!black] at (-0.75,0.62) {peso};
    \node[font=\Lato\tiny, anchor=east, text=neutral!50!black] at (-0.75,0) {bit};
  \end{tikzpicture}
  \caption{Il byte \code{00101101}: ogni bit a 1 aggiunge il peso della sua
           posizione.}
  \label{fig:binario}
\end{figure}

Per passare dal binario al decimale, quindi, si sommano i pesi delle
posizioni in cui c'è un 1. Il passaggio inverso si fa con le
\textbf{divisioni successive per 2}: si divide il numero per 2, si annota il
resto (0 o 1), si ripete sul quoziente fino ad arrivare a 0; i resti, letti
\textbf{dal basso verso l'alto}, sono le cifre binarie.

\begin{affianco}[0.36][c]
\begin{immagine}
\begin{tikzpicture}[font=\FiraCodeNL\small]
  \foreach \num/\quoz/\resto [count=\i from 0] in {13/6/1, 6/3/0, 3/1/1, 1/0/1} {
    \node[anchor=east] at (0,-\i*0.55) {\num\ : 2 = \quoz};
    \node[anchor=west, text=accent!80!black] at (0.3,-\i*0.55) {resto \textbf{\resto}};
  }
  \draw[wtarrowacc] (2.75,-1.75) -- node[right, font=\Lato\scriptsize, text=accent!70!black, align=left] {leggere\\dal basso} (2.75,0.15);
\end{tikzpicture}
\end{immagine}
\risultato
Per esempio, 13 in binario è \code{1101}: lo si verifica sommando i pesi,
$8 + 4 + 1 = 13$. Perché funziona? Ogni divisione per 2 “stacca” l'ultima
cifra binaria del numero, esattamente come dividere per 10 un numero
decimale ne stacca l'ultima cifra (il resto di $352 : 10$ è 2, l'ultima
cifra). Per questo il primo resto trovato è la cifra più a destra, e i resti
vanno letti al contrario.
\end{affianco}

\subsection{Quanti numeri in un byte}

Con $n$ posizioni e 2 simboli possibili per ciascuna si formano $2^n$
combinazioni diverse: ogni bit in più \textbf{raddoppia} le possibilità.
Un byte ne ha $2^8 = 256$, quindi rappresenta i numeri naturali da 0 a 255:

\begin{center}
\renewcommand{\arraystretch}{1.15}
\begin{tabular}{@{}>{\FiraCodeNL}c c@{\hspace{1.4cm}}>{\FiraCodeNL}c c@{}}
  \toprule
  \normalfont\textbf{bit} & \textbf{valore} & \normalfont\textbf{bit} & \textbf{valore} \\
  \midrule
  00000000 & 0 & 00000100 & 4 \\
  00000001 & 1 & 00000101 & 5 \\
  00000010 & 2 & \dots    & \dots \\
  00000011 & 3 & 11111111 & 255 \\
  \bottomrule
\end{tabular}
\end{center}

\subsection{I numeri negativi}

E i numeri negativi? Bisogna usare una parte dell'informazione per il
\textbf{segno}. L'idea più semplice, vista a lezione, è la
\term{rappresentazione in modulo e segno}: il primo bit indica il segno (0
positivo, 1 negativo) e gli altri 7 il valore assoluto. Con 7 bit si
rappresentano $2^7 = 128$ valori assoluti, da 0 a 127, quindi l'intervallo
è $[-127,\ 127]$.

\warning{Una svista nelle slide}{
  Le slide indicano l'intervallo $[-127 .. 128]$, ma in modulo e segno il
  massimo è $+127$ (\code{01111111}), non $128$: con 7 bit il valore
  assoluto più grande è $2^7 - 1 = 127$. L'intervallo è simmetrico,
  $[-127,\ 127]$, e lo zero compare due volte: \code{00000000} ($+0$) e
  \code{10000000} ($-0$).}

Proprio quel “doppio zero”, e il fatto che le somme con numeri negativi
richiederebbero circuiti speciali, fanno sì che i computer reali usino
un'altra rappresentazione, il \term{complemento a due}, in cui con un byte
si rappresentano i numeri da $-128$ a $+127$. L'idea si vede bene con 3
bit (figura~\ref{fig:ruota}): i numeri sono disposti su una ruota, come le
cifre del contachilometri di un'auto; partendo da \code{000} e andando
\textbf{indietro} di un passo si arriva a \code{111}, che quindi vale $-1$.

\begin{figure}[H]
  \centering
  \begin{tikzpicture}[font=\FiraCodeNL\scriptsize]
    \draw[neutral!40!white, line width=6pt] (0,0) circle (1.75);
    \foreach \bits/\val/\pos in {000/0/90, 001/1/45, 010/2/0, 011/3/-45,
                                  100/-4/-90, 101/-3/-135, 110/-2/180, 111/-1/135} {
      \node[mcella, minimum width=0.85cm] at (\pos:1.75) {\bits};
      \node[font=\Lato\small\bfseries, text=primary!60!black] at (\pos:2.5) {\val};
    }
    \draw[wtarrowacc, line width=1pt] (100:1.25) arc[start angle=100, end angle=125, radius=1.25];
    \node[font=\Lato\tiny, text=accent!70!black, align=center] at (0,0.3) {un passo indietro\\da 0 è $-1$};
    \node[font=\Lato\tiny, text=neutral!35!black, align=center] at (0,-0.45) {il primo bit\\è 1 per i negativi};
  \end{tikzpicture}
  \caption{Il complemento a due su 3 bit: 8 combinazioni, da $-4$ a $+3$,
           un solo zero. Su 8 bit si va da $-128$ a $+127$.}
  \label{fig:ruota}
\end{figure}

Il complemento a due ha conseguenze molto concrete in C (per esempio su
che cosa succede quando un numero “trabocca” oltre il massimo): le
vedremo nel capitolo~\ref{chap:tipi}.

\section{I tipi di dato}

\subsection{Gli stessi bit, significati diversi}

Il byte \code{01000001} contiene semplicemente otto bit. È il numero 65?
È la lettera \textit{A}? Dipende da come lo si \textbf{interpreta}: la
stessa quantità di memoria può rappresentare un numero naturale, un intero
con segno, un numero reale o un carattere. La memoria non sa che cosa
contiene; lo sa il programma.

\begin{figure}[H]
  \centering
  \begin{tikzpicture}[font=\Lato\scriptsize]
    \node[mcella, minimum width=2.6cm, minimum height=0.7cm, font=\FiraCodeNL\normalsize] (b) at (0,0) {01000001};
    \foreach \ang/\testo in {150/{come naturale:\\\textbf{65}}, 30/{come carattere:\\\textbf{A} (codice ASCII 65)},
                        -150/{come intero con segno:\\\textbf{+65}}, -30/{come parte di un \code{float}:\\dipende dagli altri 3 byte}} {
      \node[wtnodeplain, text width=3.3cm] (n) at (\ang:3.0) {\testo};
      \draw[wtarrow] (b) -- (n);
    }
  \end{tikzpicture}
  \caption{Lo stesso byte, letto con occhiali diversi.}
  \label{fig:occhiali}
\end{figure}

\subsection{Il tipo come astrazione}

I linguaggi di alto livello risolvono il problema con i \term{tipi di
dato}: si scrive il programma in termini di \textit{che cosa} sono i dati,
senza preoccuparsi di come e dove saranno davvero memorizzati. In C si può
scrivere

\begin{lstlisting}[style=c]
int variabile = 5;
\end{lstlisting}

e non ci si chiede in quali celle finirà il 5, né come sarà scritto in bit:
si usa \code{int} nel suo significato astratto di “contenitore di numeri
interi”. Eppure un \code{int} non sta in una cella sola: ne occupa
\textbf{quattro consecutive} (32 bit, quindi valori ben oltre 255). È il
tipo a stabilire che quelle quattro celle vanno lette insieme, come un
unico numero (figura~\ref{fig:int4}).

\begin{figure}[H]
  \centering
  \begin{tikzpicture}[font=\FiraCodeNL\scriptsize]
    \noindici
    \celle{c}{(0,0)}{00000101, 00000000, 00000000, 00000000}
    \foreach \i/\a in {0/2000, 1/2001, 2/2002, 3/2003}
      \node[mindir, anchor=north] at ([yshift=-1pt]c-\i.south) {\a};
    \draw[wtbrace] ([yshift=0.15cm]c-3.north east) -- ([yshift=0.15cm]c-0.north west);
    \node[font=\Lato\scriptsize, anchor=south, align=center] at ([yshift=0.45cm]$(c-0.north)!0.5!(c-3.north)$)
      {\FiraCodeNL variabile\Lato\ (\FiraCodeNL int\Lato): \textbf{4 byte} letti insieme come il numero 5};
  \end{tikzpicture}
  \caption{Un \code{int} occupa 4 celle consecutive (sui PC il byte meno
           significativo sta all'indirizzo più basso).}
  \label{fig:int4}
\end{figure}

Ogni linguaggio definisce i propri tipi. In C, i principali sono:

\begin{center}
\renewcommand{\arraystretch}{1.3}
\begin{tabular}{@{}>{\FiraCodeNL}l l l@{}}
  \toprule
  \normalfont\textbf{tipo} & \textbf{che cosa rappresenta} & \textbf{dimensione tipica} \\
  \midrule
  char   & un carattere (codice ASCII)     & 8 bit (1 byte) \\
  int    & un numero intero                & 32 bit (4 byte) \\
  float  & un numero reale                 & 32 bit (4 byte) \\
  double & un numero reale, più preciso    & 64 bit (8 byte) \\
  \bottomrule
\end{tabular}
\end{center}

Le dimensioni sono quelle dei computer di oggi; lo standard del C fissa solo
dei minimi, come vedremo nel capitolo~\ref{chap:tipi}.

\section{La gerarchia di memoria}

La memoria centrale può contenere molti programmi e dati, ma ha due difetti:

\begin{itemize}
  \item è \textbf{lenta} rispetto al processore: leggere una cella richiede
        il tempo di un centinaio di operazioni della CPU;
  \item è \textbf{volatile}: quando si spegne il computer, il contenuto si
        perde.
\end{itemize}

Contro la lentezza si costruisce una \term{gerarchia}: piccole memorie
velocissime vicino al processore, memorie via via più grandi e lente man
mano che ci si allontana. Le memorie veloci costano molto per ogni bit,
quindi sono piccole. Contro la volatilità ci sono le \term{memorie di massa}
(dischi, SSD, chiavette), che conservano i dati anche a computer spento
(figura~\ref{fig:gerarchia}).

\begin{figure}[H]
  \centering
  \begin{tikzpicture}[font=\Lato\scriptsize]
    \def\h{0.78}
    \foreach \n/\nome/\vel/\dim/\col [count=\i from 0] in {
        0/registri (dentro la CPU)/circa 1 ns/centinaia di byte/accent,
        1/cache/circa 10 ns/kilobyte\,--\,megabyte/accent,
        2/memoria centrale (RAM)/circa 100 ns/gigabyte/primary,
        3/memoria di massa (disco{,} SSD)/da 0{,}1 a 10 ms/terabyte/green-gradient-6,
        4/archivio (nastri{,} cloud)/secondi/molti terabyte/neutral} {
      \pgfmathsetmacro{\w}{2.9+\i*1.3}
      \fill[\col!14!white, draw=\col!65!black, line width=0.7pt]
        (-\w/2,-\i*\h) -- (\w/2,-\i*\h) -- (\w/2+0.65,-\i*\h-\h) -- (-\w/2-0.65,-\i*\h-\h) -- cycle;
      \node[font=\Lato\scriptsize\bfseries] at (0,-\i*\h-\h/2) {\nome};
      \node[anchor=west] at (5.2,-\i*\h-\h/2) {\vel};
      \node[anchor=west] at (7.6,-\i*\h-\h/2) {\dim};
    }
    \node[anchor=west, font=\Lato\scriptsize\bfseries] at (5.2,0.3) {tempo di accesso};
    \node[anchor=west, font=\Lato\scriptsize\bfseries] at (7.6,0.3) {dimensione};
    \draw[wtarrow, <->] (-5.3,-0.1) -- node[wtlab, rotate=90, fill=white] {più veloci, più costose} (-5.3,-3.8);
    \draw[neutral!50!white] (10.2,-0.25) -- (10.2,-2.25);
    \node[font=\Lato\tiny, text=neutral!35!black, anchor=west] at (10.3,-1.25) {volatili};
    \draw[neutral!50!white] (10.2,-2.45) -- (10.2,-3.85);
    \node[font=\Lato\tiny, text=neutral!35!black, anchor=west] at (10.3,-3.15) {permanenti};
  \end{tikzpicture}
  \caption{La gerarchia di memoria (1 ns è un miliardesimo di secondo). In
           alto memorie piccole, veloci, costose e volatili; in basso grandi,
           lente, economiche e permanenti.}
  \label{fig:gerarchia}
\end{figure}

Un'immagine aiuta: la CPU è un impiegato alla scrivania. I registri sono ciò
che tiene in mano; la cache, i fogli sulla scrivania; la RAM, lo scaffale
nella stanza; il disco, l'archivio in cantina. L'impiegato lavora solo su
ciò che ha in mano, e se un foglio non è sulla scrivania deve alzarsi a
prenderlo: per questo i dati usati spesso vengono tenuti vicini.

\section{Il software}

\subsection{Gli strati del software}

Fra l'hardware e i programmi che usiamo c'è un sistema software a strati
(figura~\ref{fig:strati}). Ogni strato usa i servizi di quello sotto e ne
offre di più comodi a quello sopra: è ancora l'idea di astrazione.

\begin{figure}[H]
  \centering
  \begin{tikzpicture}[font=\Lato\scriptsize,
      strato/.style={draw=#1!65!black, fill=#1!12!white, line width=0.7pt,
        minimum width=9.5cm, minimum height=0.68cm, anchor=north},
      es/.style={anchor=west, font=\Lato\scriptsize, text=neutral!25!black}]
    \node[strato=accent] (app) at (0,0) {\textbf{applicazioni}};
    \node[strato=primary] (lib) at ([yshift=-2pt]app.south) {\textbf{librerie} (per esempio la libreria standard del C)};
    \node[strato=green-gradient-6] (int) at ([yshift=-2pt]lib.south) {\textbf{interfaccia}: a riga di comando (la shell) o grafica};
    \node[strato=orange-gradient-6] (so) at ([yshift=-2pt]int.south) {\textbf{sistema operativo}: gestione di processi, memoria, periferiche, file};
    \node[strato=red-gradient-6] (ker) at ([yshift=-2pt]so.south) {\textbf{nucleo} (\textit{kernel}) del sistema operativo};
    \node[strato=neutral] (hw) at ([yshift=-2pt]ker.south) {\textbf{hardware}: CPU, memoria, dispositivi};
    \node[es] at ([xshift=0.4cm]app.east) {browser, VS Code, i nostri programmi};
    \node[es] at ([xshift=0.4cm]lib.east) {\code{printf}, \code{scanf}\dots};
    \node[es] at ([xshift=0.4cm]int.east) {bash, il desktop};
    \node[es] at ([xshift=0.4cm]so.east) {Linux, Windows, macOS};
  \end{tikzpicture}
  \caption{Fra l'hardware e l'ambiente in cui si sviluppano i programmi
           c'è un sistema software stratificato.}
  \label{fig:strati}
\end{figure}

\subsection{L'avvio del computer}

Il primo programma eseguito all'accensione è il \term{sistema operativo}:
è lui che, quando chiediamo di eseguire un programma, lo carica in memoria
e dice al processore dove si trova la sua prima istruzione. Ma chi dice al
processore dove si trova la prima istruzione del sistema operativo? È un
problema dell'uovo e della gallina, risolto così
(figura~\ref{fig:avvio}): all'accensione il computer esegue un piccolo
programma che sta in una memoria permanente della scheda madre (il
\textit{firmware}); questo carica in memoria il contenuto di un blocco
predefinito di un disco, il \term{Master Boot Record} (MBR). Il codice del
MBR, chiamato \term{bootloader}, è una sequenza di byte che si assume essere
un programma in linguaggio macchina: il suo compito è caricare il sistema
operativo. Il MBR è configurabile: se sul disco ci sono più sistemi
operativi, si può scegliere quale avviare.

\begin{figure}[H]
  \centering
  \begin{tikzpicture}[font=\Lato\scriptsize,
      passo/.style={wtnode, font=\Lato\scriptsize, text width=2.8cm, minimum height=1.3cm}]
    \node[passo, draw=neutral!60!white, fill=neutral!8!white, text width=2.4cm] (a) at (0,0) {\textbf{accensione}};
    \node[passo] (b) at (3.6,0) {\textbf{firmware}\\(BIOS o UEFI)\\nella scheda madre};
    \node[passo] (c) at (7.4,0) {\textbf{bootloader}\\letto dal MBR del disco};
    \node[passo, draw=accent!60!black, fill=accent!8!white] (d) at (11.2,0) {\textbf{sistema operativo}\\caricato in memoria};
    \node[passo, draw=green-gradient-6, fill=green-gradient-1!45!white, text width=2.4cm] (e) at (14.8,0) {i \textbf{programmi} dell'utente};
    \foreach \x/\y in {a/b, b/c, c/d, d/e} \draw[wtarrow] (\x) -- (\y);
  \end{tikzpicture}
  \caption{L'avvio: ogni programma carica il successivo, più grande e più
           capace.}
  \label{fig:avvio}
\end{figure}

\info{MBR e UEFI}{
  Il Master Boot Record è il meccanismo storico dei PC. I computer recenti
  usano il firmware \textit{UEFI}, che invece di un singolo blocco legge il
  bootloader da una piccola partizione dedicata del disco. L'idea non
  cambia: un programma minuscolo, che sta in un posto noto in anticipo,
  avvia un programma più grande.}

\subsection{Che cosa fa il sistema operativo}

Il sistema operativo \textbf{alza il livello di astrazione} della macchina:
è più facile da programmare dell'hardware sottostante e semplifica la
scrittura di ogni altro programma. Le sue funzioni principali:

\begin{itemize}
  \item \textbf{gestione dei processi}, cioè dei programmi in esecuzione:
        li crea e li termina, li sospende e li riattiva, li fa comunicare.
        Il processore esegue un'istruzione alla volta, ma il sistema
        operativo lo assegna a turno ai vari programmi per pochi millesimi
        di secondo ciascuno (la \textit{schedulazione}): così sembrano
        procedere tutti insieme;
  \item \textbf{gestione della memoria centrale}: assegna la memoria ai
        processi, li carica e li rimuove, e la \textbf{protegge}, cioè
        impedisce a un programma di scrivere nella memoria di un altro;
  \item \textbf{gestione delle periferiche}: le configura e le pilota, con
        il \textit{buffering} (tiene i dati in memoria finché non vengono
        elaborati) e il \textit{caching} (tiene vicini i dati usati spesso);
  \item \textbf{gestione del file system}: nasconde all'utente come sono
        organizzati i dati sul disco, offrendo il concetto di \textit{file}
        e di \textit{cartella} (l'albero di cartelle del
        capitolo~\ref{chap:ambiente}), con le operazioni di creazione,
        lettura, scrittura, copia, ricerca e i diritti di accesso;
  \item \textbf{interfaccia utente}: a riga di comando (l'interprete dei
        comandi, la shell che usiamo nel terminale) o grafica;
  \item \textbf{rilevamento degli errori}: divisione per zero, carta
        finita, rete caduta\dots; per ogni errore serve un'azione di
        recupero appropriata.
\end{itemize}

L'ultima voce ci riguarderà da vicino. Un programma C che divide un intero
per zero chiede al processore un'operazione impossibile; il processore
avvisa il sistema operativo, che \textbf{termina il programma}:

\begin{affianco}[0.45][c]
\begin{lstlisting}[style=c, name=div0.c]
#include <stdio.h>

int main(void)
{
    int a = 7, b = 0;
    printf("%d\n", a / b);
    return 0;
}
\end{lstlisting}
\risultato
\begin{console}
$ gcc -o div0 div0.c
$ ./div0
Floating point exception
\end{console}
\wtnota{Il messaggio parla di “virgola mobile” anche se la divisione è
fra interi: è un nome storico.}
\end{affianco}

\subsection{I servizi per i programmi}

Il sistema operativo offre servizi ai programmi: caricarli in memoria ed
eseguirli, compiere le operazioni di I/O (un programma normale non può
pilotare direttamente i dispositivi), leggere e scrivere i file, far
comunicare processi sullo stesso computer o su computer diversi (con memoria
condivisa o scambio di messaggi), rilevare gli errori. Ci sono poi servizi
per gestire il sistema nel suo insieme: allocazione delle risorse fra utenti
e processi, contabilità del loro uso, protezione e sicurezza.

Anche il più semplice dei nostri programmi passa da questi servizi. Quando
il programma del capitolo precedente esegue \code{printf}, il testo non va
direttamente sullo schermo (figura~\ref{fig:printf-so}).

\begin{figure}[H]
  \centering
  \begin{tikzpicture}[font=\Lato\scriptsize,
      p/.style={wtnode, font=\Lato\scriptsize, text width=3.5cm, minimum height=1.15cm}]
    \node[p] (a) at (0,0) {il nostro programma chiama \code{printf}};
    \node[p] (b) at (4.5,0) {la \textbf{libreria} del C prepara il testo};
    \node[p, draw=orange-gradient-6!80!black, fill=orange-gradient-1!30!white] (c) at (9.0,0) {chiede al \textbf{sistema operativo} di scriverlo};
    \node[p, draw=green-gradient-6, fill=green-gradient-1!45!white] (d) at (13.5,0) {il testo compare nel \textbf{terminale}};
    \foreach \x/\y in {a/b, b/c, c/d} \draw[wtarrow] (\x) -- (\y);
  \end{tikzpicture}
  \caption{Il percorso di una \code{printf}: dal programma alla libreria,
           dalla libreria al sistema operativo, e infine al dispositivo.}
  \label{fig:printf-so}
\end{figure}

\subsection{Installare e scrivere un sistema operativo}

Lo stesso sistema operativo gira su macchine diverse, quindi va
\textbf{configurato} in fase di installazione: processore, memoria,
suddivisione del disco (partizioni), dispositivi presenti, opzioni come
l'algoritmo di schedulazione o il numero massimo di processi.

I primi sistemi operativi erano scritti in assembly; oggi si scrivono in
linguaggi di alto livello, perché il codice si scrive più in fretta, è più
compatto, più facile da capire e da correggere. Un sistema operativo scritto
in un linguaggio di alto livello è molto più flessibile e manutenibile. Non
a caso, il C nacque proprio per scrivere un sistema operativo, Unix, come
racconta il capitolo~\ref{chap:strutturata}.

\insintesi{
  Un calcolatore è fatto di CPU, memoria e dispositivi di I/O, collegati
  dal bus. La memoria è una fila di celle di un byte, ognuna con un
  indirizzo, che contengono sia i dati sia le istruzioni dei programmi. La
  CPU sa eseguire solo istruzioni elementari del proprio
  \textbf{linguaggio macchina}, una alla volta, guidata dal program counter;
  ripetizioni e scelte si ottengono con i salti. I linguaggi di alto
  livello come il C nascondono questi dettagli (astrazione), ma vanno
  \textbf{tradotti}. I bit di per sé non hanno significato: sono i
  \textbf{tipi} a stabilire come interpretarli. Il sistema operativo, il
  primo programma caricato all'avvio, gestisce processi, memoria, periferiche
  e file, e offre i suoi servizi a tutti gli altri programmi.}

\section{Riferimenti}

\begin{itemize}
  \item \textbf{Deitel \& Deitel}, \textit{Il linguaggio C --- Fondamenti e
        tecniche di programmazione}, Pearson: capitolo~1 (computer,
        hardware e software, linguaggi macchina, assembly e di alto
        livello).
  \item \textbf{Treccani}, voce \textit{programmazione}. \\
        \urlx{https://www.treccani.it/vocabolario/programmazione/}
\end{itemize}
